% paper/sections/09f_pressure_summary.tex
%
% §9 章末まとめ：圧力ジャンプ型 FCCD-PPE の整理
% ═══════════════════════════════════════════════════════════════════════════════
\subsection{PPE 解法の精度整理：本質問題と閉包条件}
\label{sec:pressure_accuracy_summary}
% ═══════════════════════════════════════════════════════════════════════════════

本章の結論は，複数の圧力解法を状況に応じて同列に併用することではない．
二相の圧力射影段階の本質問題を
\[
\begin{gathered}
  \text{密度ジャンプ}
  +\text{圧力ジャンプ}
  +\text{面演算子整合}
  +\text{高次片側データ},\\
  \text{表面エネルギー毛管閉包}
  +\text{平衡判定}
  +\text{欠陥補正},\\
  \text{壁付き face-state 制約}
  +\text{大域ゲージ・再投影分離}
  +\text{診断契約}
\end{gathered}
\]
に分解し，各問題に必要な手段だけを残すことである．
この整理は，射影法の圧力補正・境界条件文献
~\cite{Chorin1968,vanKan1986,BrownCortezMinion2001,Guermond2006,PyoShen2007}
と，界面ジャンプ・CSF/BF 文献
~\cite{Brackbill1992,Fedkiw1999,Kang2000,Francois2006,Popinet2009}
の役割を分解して読み直すものである．

\begin{table}[htbp]
\centering
\caption{§9 で扱った圧力閉包対象と閉包条件の整理}
\label{tab:accuracy_summary}
\PaperTableDenseSetup
\begin{tabularx}{\linewidth}{@{}>{\raggedright\arraybackslash}p{0.18\linewidth}>{\raggedright\arraybackslash}X>{\raggedright\arraybackslash}X>{\raggedright\arraybackslash}p{0.18\linewidth}@{}}
\toprule
閉包対象 & 破綻機構 & 本稿の閉包条件 & 役割を限定するモデル \\
\midrule
密度ジャンプ & smoothed-Heaviside 一括 PPE が $\rho^{-1}$ の不連続を微分し，条件数と界面誤差が増幅する &
相内定数密度 Poisson &
smoothed-Heaviside 一括 PPE は低密度比の対照モデル \\
圧力ジャンプ & Young--Laplace 圧力差を連続圧力として扱うと界面応力が曖昧になる &
$j_{gl}=p_g-p_l=-\sigma\kappa_{lg}$ を明示する向き付き条件 &
符号を持たない $\sigma\kappa_{lg}$ 場ではなく，向き付き圧力差として扱う \\
面演算子整合 & CSF 体積力と圧力勾配が別の離散位置・別の演算子に乗る &
$G_\Gamma(p;j)=G(p)-B(j)$ のジャンプ補正付き面勾配 &
CSF は Balanced--Force 原理を説明する対照モデル \\
圧力反力代表 & スカラー PPE の発散だけを合わせると，圧力仕事に直交しない発散ゼロ代表差 $Z$ が残る &
$G_A=-M_A^{-1}D_f^TW_p$ と $L_A=D_fG_A$ を同じ面計量で組にする &
raw compact 勾配は，この随伴性が証明されない限り本番圧力反力ではない \\
補助場相殺 & $p_\mathrm{comp}=\tilde p+J_\psi$ では楕円求解が $J_\psi$ を吸収できる &
ジャンプ補正付き面勾配を本章の閉包にする &
滑らかな補助場による圧力合成は主閉包にしない \\
高次片側データ & CCD/FCCD ステンシルが不連続圧力を跨ぐと設計次数が崩れる &
HFE によるジャンプ補正済み Hermite 延長 &
界面を跨いだ通常 CCD 強行は不可 \\
毛管補正量 & 表面エネルギーの離散第一変分と PPE・補正段の面空間がずれると，静的液滴には寄生流れを，非平衡界面には過小駆動を生む &
表面エネルギーの Riesz 表現 $s_f$，成分体積反力 $B_h\mu$，成分補正済み補正量 $\widehat c_{\sigma,f}=s_f-B_h\mu$ &
曲率上限・人工減衰・時間刻み縮小による症状処理ではない \\
壁付き face-state & 圧力射影後に壁速度だけを修復すると，輸送 face state と公開 nodal state が分裂する &
$F_w=\ker C_w$ を速度状態空間とし，$P_w$ と $G_w=P_wG_A$ を制約付き空間の診断 gate とする &
post-pressure wall-only projection と full-space additive KKT は本番閉包にしない \\
診断契約 & legacy 節点曲率や部分 PPE 右辺を active operator の量として読むと誤判定する &
active face curvature，solver 直前 RHS，$M_A=Q_f/\alpha_f$ 上の制約付き Hodge residual を記録する &
\texttt{kappa\_max} や CSF 型 residual を pressure-jump 本番力の代用判定にしない \\
高次演算子処理 & 高次評価演算子と低次補正演算子を同一視すると，設計次数と補正安定性の役割分担が崩れる &
DC 残差受理条件 $\eta_{\mathrm{DC}}\le\varepsilon_{\mathrm{PPE}}$: 高次評価 $L_H$ + 低次補正 $L_L$ &
直接解型 CCD は滑らかな製造解・成分検証に限定 \\
零モード・再投影 & ジャンプ補正接続演算子の全体定数零モードと格子再生成後の再投影 PPE への圧力ジャンプ持ち込みが人工エネルギーを生む &
大域ゲージと再投影 PPE での界面ジャンプ情報の分離 &
直前時刻 $j_{gl}$ の再投影継承は禁止 \\
\bottomrule
\end{tabularx}
\end{table}

\begin{tcolorbox}[mybox, title={閉包条件：§9 の圧力ジャンプ型閉包}]
\label{box:ppe_design_rationale}
§9 の圧力ジャンプ型閉包は次の同時成立を要求する：
\[
  \boxed{
  \begin{gathered}
  \text{相内定数密度 PPE}
  +G_A=-M_A^{-1}D_f^TW_p
  +L_A=D_fG_A \\
  +G_\Gamma(p;j)=G_A(p)-B(j)
  +\mathrm{HFE} \\
  +\widehat c_{\sigma,f}\ \text{（表面エネルギー・成分体積制約付き鞍点系）} \\
  +F_w=\ker C_w,\quad G_w=P_wG_A\ \text{（壁付き face-state 診断 gate）} \\
  +\mathrm{DC}\{\eta_{\mathrm{DC}}\le\varepsilon_{\mathrm{PPE}}\}
  +\text{大域ゲージ・再投影分離・active 診断契約}
  \end{gathered}
  }.
\]
ここでジャンプ補正付き面勾配は毛管波専用の修正ではない．
液滴，気泡，毛管波，平坦 Rayleigh--Taylor 界面のいずれでも，
同じ $j_{gl}=p_g-p_l=-\sigma\kappa_{lg}$ と
同じ $G_\Gamma(p;j)$ を使うことが閉包条件である．
非一様格子ではこの受理条件を
$G_{\Gamma,f}^{\mathrm{nu}}=G_f^{\mathrm{nu}}-B_f^{\mathrm{nu}}$，
$B_f^{\mathrm{nu}}=s_fj_f/H_f$ と読み替え，
PPE 右辺と速度補正が同じ $H_f$・同じ面係数・同じ非一様発散を共有する．
さらに毛管補正量は再初期化前の輸送終点上の
表面エネルギーの仮想仕事から $s_f$ として構成し，
同じ圧力仕事と随伴な面計量上の成分体積制約付き鞍点系で
$\widehat c_{\sigma,f}=s_f-B_h\mu$ に補正してから速度補正へ渡す．
このとき $G_A$ はスカラー PPE のための任意の高次勾配ではなく，
式~\eqref{eq:pressure_reaction_green_identity} を満たす圧力反力代表である．
したがって $D_f\widetilde G=D_fG_A$ を満たすだけの
発散同値な代表 $\widetilde G$ は，圧力仕事と毛管 Hodge 分解の意味では
同じスキームではない．
式~\eqref{eq:capillary_range_projection} の $L_A/Z_A$ は
静的平衡判定と原因診断のための代表であり，標準毛管力を置換しない．
壁付き計算では，この判定を full face 空間ではなく
式~\eqref{eq:wall_constrained_face_space} の $F_w$ と
active pressure metric $M_A=Q_f/\alpha_f$ 上で読む．
圧力射影後に壁トレースだけを修復する方法，full-space additive KKT を
そのまま本番閉包とする方法，および $s_f\leftarrow L_A(s_f)$ による
blanket range replacement は採用しない．

この閉包が保証するのは，PPE と速度補正が同じ面フラックス整合を共有し，
圧力ジャンプが滑らかな圧力自由度として消えないことである．
また，圧力ジャンプ経路で物理的に仕事をする圧力は
$A_f(G_fp-B_f(j))$ の面上加速度であり，
拡散界面帯内の未補正節点圧力は相内スカラー圧力そのものではない．
したがってスナップショットのスカラー圧力は保存された面上加速度から復元した
同相内圧力代表として定義し，未補正節点値は診断用の内部自由度として扱う．
界面帯を ``未定義'' として隠す表示や未補正節点値を相内圧力として読む表示は，
本章の圧力代表ではない．
一方で，曲率上限，HFE の実効精度，毛管 CFL，壁面影響，長時間の幾何エネルギー安定性は
本節の保証範囲外であり，後続の統合検証で評価する．
\end{tcolorbox}

\begin{tcolorbox}[warnbox, title={§9 が主張しないもの}]
\begin{itemize}[noitemsep]
  \item \textbf{smoothed-Heaviside 一括 PPE}：
        低・中密度比の対照モデルとして扱うが，高密度比の圧力ジャンプ型閉包の解決策とはしない．
  \item \textbf{CSF 体積力}：
        Balanced--Force 原理を説明する基準モデルとして扱うが，最終的な毛管閉包とはしない．
  \item \textbf{滑らかな補助場による圧力合成}：
        既知ジャンプを自由圧力へ吸収し得るため，圧力ジャンプ型設定の既定にはしない．
  \item \textbf{有限体積法/有限差分低次補正}：
        DC の低次補正演算子や対照基準として使うが，CCD/FCCD 面演算子の代替主張にはしない．
  \item \textbf{post-pressure wall-only repair}：
        壁トレースを消しても発散と face/nodal state 同一性を同時に保つとは限らないため，
        壁付き標準経路では $F_w=\ker C_w$ を状態空間として先に定義する．
  \item \textbf{legacy diagnostic の物理量化}：
        節点曲率最大値，途中段階の PPE 右辺，CSF 型 balanced-force residual を，
        pressure-jump 経路の active 毛管力・active 右辺・active Hodge 残差として読まない．
\end{itemize}
\end{tcolorbox}

したがって，9 章を読んだ後に残る物語は単純である：
圧力射影段階の本質は「高次 Poisson 求解」ではなく，
不連続な界面応力を同じ面演算子上で保存したまま投影することである．
そのために，分相化で密度ジャンプを外し，
向き付き界面条件で Young--Laplace ジャンプを PPE に入れ，
HFE で高次ステンシルの滑らかさを守り，
毛管補正量を表面エネルギー仮想仕事と成分体積制約付き鞍点系で閉じ，
DC で高次演算子を扱い，
大域ゲージと再投影分離で閉じる．
